\section{Introduction}
\label{sec:introduction}

Near a polystable bundle, a finite-dimensional deformation
space describes the loss of rigidity of Hermitian--Einstein
metrics. When the underlying Kähler manifold also degenerates,
this description requires estimates uniform in the geometric
parameters. An additive error in a finite operator is
insufficient if some eigenvalues are much smaller than
others. Likewise, a stationary reduction alone does not
determine the trajectory selected by a spatially varying
initial metric or the harmonic metric on nearby bundle
classes. This paper proves these three comparisons on one
nonflat bundle family over compact conifold smoothings.

\subsection{The geometric family and the principal theorem}

We use a polarized family $X_p$ of Calabi--Yau threefolds
whose central fibre has eight ordinary double points.
The base is a small ball in the six-dimensional space
\begin{equation}\label{eq:intro-base}
 \mathcal R=\left\{p\in\CC^8:
       \sum_{i\in I_f}p_i=\sum_{i\in I_g}p_i=0\right\},
 \qquad |I_f|=|I_g|=4.
\end{equation}
Its smooth locus is the complement of the eight coordinate
hyperplanes. The node parameters satisfy
$\tau_i(p)=p_i/u_i(p)$ for nonvanishing holomorphic units.
The smooth fibres are complete intersections of bidegrees
$(1,1),(1,1),(2,2)$ in $\PP^3\times\PP^3$.
Write $H=\OO(2,2)$ and $L=\OO(1,-1)$ on each fibre.
The actual Ricci-flat metric in $c_1(H)$ has volume
$V=160/3$, and the split bundle is
\[
                 E^{\rm sp}=L\otimes\CC^2\oplus L^{-1}\otimes\CC^2.
\]
The polarized construction and the uniform actual-metric
input are restated in Section~\ref{sec:geometry} from
\cite{W_Hodge}. All subsequent bundle and flow estimates
are proved here.

The bundle deformation parameters are two matrices
$z=(A,B)\in W=\operatorname{Mat}_2(\CC)\oplus\operatorname{Mat}_2(\CC)$,
with compact action
$(k_+,k_-)\cdot(A,B)=(k_+Ak_-^{-1},k_-Bk_+^{-1})$.
The two harmonic extension forms have squared norms
$c_+(p),c_-(p)>0$. They are uniformly bounded above
and below and converge to positive central values.
On the real seven-dimensional space
\[
 \mathfrak h=\{(S_+,S_-):S_\pm=S_\pm^*,\quad\tr S_++\tr S_-=0\}
\]
define the mass and energy
\begin{equation}\label{eq:intro-finite-pair}
 \begin{split}
 M(S)&=V(\tr S_+^2+\tr S_-^2),\\
 E_{p,z}(S)&=2c_+\|S_+A-AS_-\|_{\rm HS}^2
                  +2c_-\|S_-B-BS_+\|_{\rm HS}^2.
 \end{split}
\end{equation}
Let $\mu_1\le\cdots\le\mu_7$ be its generalized
eigenvalues, including zeros. The quadratic moment is
\[
 \mu_{2,p}(A,B)=V^{-1}
       (c_+AA^*-c_-B^*B,\ c_-BB^*-c_+A^*A).
\]
Its zero set becomes $XX^*=Y^*Y$, $YY^*=X^*X$
under $X=\sqrt{c_+}A$, $Y=\sqrt{c_-}B$.

\begin{theorem}\label{thm:main}
On the family above there is an actual holomorphic
rank-four bundle family $\mathcal E_{p,w}$ realizing
the full local deformation germ of $E^{\rm sp}$, with
$\det\mathcal E_{p,w}=\OO$ and $c_2(\mathcal E_{p,w})H=32$.
After one common shrink, the following hold uniformly
over every smooth $p$ and all small bundle parameters.

\smallskip\noindent
\textup{(i)} There is a real $K$-equivariant map
$\chi_p(z)=z+O(|z|^2)$, $d\chi_p=I+O(|z|)$, such that
every $z\in\mu_{2,p}^{-1}(0)$ gives an actual HE
connection on $\mathcal E_{p,\chi_p(z)}$.
These representatives cover the sufficiently small
polystable objects in the stated local germ. The germ
contains no stable rank-four objects.

\smallskip\noindent
\textup{(ii)} At these HE connections the positive
operator linearizing $h^{-1}\dot h=-2\mathrm i\Lambda F_h$
has complete small spectrum
\[
 (1-C|z|)\mu_j\le\lambda_j\le(1+C|z|)\mu_j
       \quad(1\le j\le7),\qquad\lambda_8\ge c>0.
\]
Its exact trace-free kernel is the compact infinitesimal
stabilizer. The comparison allows changing kernel dimension
and arbitrary ratios of positive eigenvalues.

\smallskip\noindent
\textup{(iii)} Fix an HE member, write $r=|\chi_p(z)|$,
and fix any finite graph regularity order $m\ge4$.
There is a common spatial $C^1$ basin of smooth
determinant-normalized initial metrics and a chosen
$C^m$ seven-dimensional invariant graph. Every flow
in that basin selects a $C^1$ finite initial phase
and tracks the corresponding exact graph trajectory
with error $CD_*e^{-ct}$ in the logarithmic $C^1$ norm.
Here $D_*$ is its initial displacement from the graph.
The selected graph trajectory remains in the actual
equation for all positive time. Its finite metric is
the pulled-back $L^2$ metric, its energy is twice
Donaldson's functional, and its Hessian at the HE
origin is relatively $1+O(|z|)$ equivalent to $E_{p,z}$
in the same $S$ coordinates. The preparation time
$4\sigma^{-1}\log(1/r)$ is negligible compared with
every positive inverse weak rate.

\smallskip\noindent
\textup{(iv)} The local HE classes are parametrized by
$\CC^2/\mathsf W$, where $\mathsf W=(\ZZ/2)^2\rtimes S_2$
acts by signs and permutation. On every quotient stratum,
their actual harmonic connection metric is relatively
$1+O(|u|)$ equivalent to $4|\dd u|^2$.
For paths confined to a fixed sufficiently small local
locus, its metric completion is the corresponding
closed parameter locus, and its unique tangent cone at
the split bundle is $(\CC^2,4|\dd u|^2)/\mathsf W$.
The rescaled distance error is $O(\rho)$ on bounded
balls, uniformly in the degenerating base.
\end{theorem}

The theorem combines the precise statements proved below.
The bundle realization is Proposition~\ref{prop:effective-bundles}.
The relative spectrum and exact balancing are
Theorem~\ref{thm:complete-bundle-spectrum} and
Proposition~\ref{prop:exact-balancing-coordinates};
Section~\ref{sec:balancing} proves holomorphic
polystability and local completeness. The dynamical
claims are Theorems~\ref{thm:spatial-entry},
\ref{thm:invariant-graph}, \ref{thm:all-time-tracking}
and Corollary~\ref{cor:relative-finite-hessian}.
The actual quotient identification, harmonic metric and
distance are Theorems~\ref{thm:unitary-rigidity},
\ref{thm:harmonic-quotient} and \ref{thm:metric-cone}.
Their proofs use the same connections and the same
compact action. The flow parameters are fixed along
each trajectory; uniformity permits them to approach
their limits jointly when comparing different trajectories.

The balanced normal form has two cycle parameters.
Its four types have trace-free stabilizer dimensions
$1,3,2,7$: two distinct nonzero cycles, a repeated
nonzero cycle, one nonzero cycle and two zero vertex
factors, and the split bundle. All occur in the theorem.
The action in \eqref{eq:intro-finite-pair} controls
both the zero rates and the unitary quotient directions.
The graph metric, however, varies metrics on a fixed
bundle, whereas the harmonic metric varies its
holomorphic isomorphism class.

A simple path exhibits two widely separated rates.
Take $X=Y=\diag(h,h+h^b)$ with fixed $b>1$.
There is one exact zero, two rates
$2h^{2b}V^{-1}(1+O(h))$, and four rates
$8h^2V^{-1}(1+o(1))$; see
Proposition~\ref{prop:two-cycle-spectrum}.
This permits an arbitrarily long weak time scale
independently of every node scale. The all-time
comparison uses the exact finite field, not a
truncation whose error could accumulate on that scale.

\subsection{Relation to earlier work}

Local HE reduction and geometric invariant theory on a
fixed compact Kähler manifold are established tools.
Buchdahl--Schumacher construct equivariant local
deformations and analytic GIT quotients
\cite[Theorems 5--8]{BS}; their Proposition~5 identifies
the harmonic form. Toda describes polystable sheaf
germs by Ext-quivers with convergent relations
\cite[Theorem 3.2]{Toda}. Delloque proves complementary
correction and local quotient results for elliptic
polynomial curvature equations, including HYM
\cite[Propositions 4.3 and 4.5 and Theorem 7.4]{Delloque}.
Here the relations vanish for a geometrically realized
family, although its threefold Ext-two group is nonzero.
The analytic estimates must also remain uniform as the
underlying metric degenerates. The dense-stable
hypothesis in \cite[Theorem 10]{BS} is absent in our
rank-four neighbourhood, so that extension theorem
does not provide the harmonic distance result used here.

Nonflat HYM metrics through conifold transitions were
constructed by Chuan \cite{Chuan}. Collins--Picard--Yau
prove a tangent-bundle theorem with quantitative local
profiles for the Fu--Li--Yau balanced backgrounds
\cite[Theorem 1.1]{CPY}. Our background is the actual
Ricci-flat Kähler metric, the bundle has rank four,
and the base and bundle parameters may have independent
scales. The additional conclusions concern the complete
relative endomorphism spectrum, the spatial metric
flow and the harmonic metric on actual bundle classes.

Haiden--Katzarkov--Kontsevich--Pandit identify iterated
weight filtrations and iterated-logarithmic asymptotics
for finite algebra flows \cite{HKKPweights}, and prove
actual bundle-flow asymptotics on curves
\cite[Theorem 1.1]{HKKPflows}. Their order and
integrable-forcing comparison in Section~4.2 of
\cite{HKKPflows} applies in arbitrary compact Kähler
dimension. Lam proves logarithmic and iterated-logarithmic
asymptotics on a fixed smooth compact Kähler base for
semistable bundles with locally free graded object and
the stated rationality condition on the invariant inner
product \cite[Theorem 8.2]{Lam}. His Proposition~8.6
compares actual and reduced trajectories at bounded
distance after a preliminary complex-gauge change.
Our theorem instead concerns a uniform small basin
around polystable HE metrics over degenerating geometry
and gives decaying spatial tracking to a selected exact
graph trajectory. It does not determine nonlinear
decay exponents or leading coefficients of a trajectory.

Invariant manifolds for parabolic equations are classical;
see, for example, \cite[Theorem 2.1]{Simonett}.
Here we prove the required $C^0$-to-$C^1$ heat estimates
from scalar Ricci-flat heat bounds and controlled split
connection coefficients. The logarithmic metric equation
depends only on its first spatial jet. These facts give
a direct graph construction without assuming uniform
high-order coordinate estimates at the shrinking nodes.
The metric-distance inequality then proves confinement
for an open spatial basin and keeps the selected graph
trajectory in the original equation.

Earlier work in this series treats stationary local
reduction on a fixed smooth background \cite{W_Local},
finite phase reduction for trajectories satisfying its
confinement hypotheses \cite{W_Phase}, and all-time
approximation on degenerating curves with the stated
spatial initial-data class \cite{W_Curves}.
The present proof verifies the complementary estimates,
confinement and quotient geometry anew on the nonflat
conifold family. Its contribution is this geometric
realization and the simultaneous uniformity of the
spectral, dynamical and harmonic-metric comparisons.

\subsection{Structure of the proof}

Sections~\ref{sec:geometry}--\ref{sec:projected-he}
construct the bundles and the uniform analytic estimates.
Sections~\ref{sec:relative-spectrum}--\ref{sec:moving-scales}
prove the complete relative spectrum, exact balancing,
polystability classification and multiscale example.
Sections~\ref{sec:spatial-entry}--\ref{sec:phase-tracking}
construct the spatial basin, invariant graph and phase.
Sections~\ref{sec:harmonic-quotient}--\ref{sec:metric-cone}
identify actual classes and prove the harmonic metric
and distance comparisons. All constants refer to the
common small domains specified in these statements.
